\begin{thebibliography}{99}
\addcontentsline{toc}{chapter}{参考资料}
\bibitem{openpi} Physical Intelligence. OpenPI项目与LIBERO运行示例. \url{https://github.com/Physical-Intelligence/openpi}. 访问日期：2026年11月15日.
\bibitem{libero} LIBERO项目. 语言条件机器人操作基准. \url{https://github.com/Lifelong-Robot-Learning/LIBERO}. 访问日期：2026年11月15日.
\bibitem{lerobot} Hugging Face. LeRobot机器人学习工具. \url{https://github.com/huggingface/lerobot}. 访问日期：2026年11月15日.
\bibitem{qwen} Qwen团队. Qwen3-VL项目. \url{https://github.com/QwenLM/Qwen3-VL}. 访问日期：2026年11月15日.
\bibitem{moveit} MoveIt项目. 控制器配置说明. \url{https://moveit.picknik.ai/humble/doc/examples/controller_configuration/controller_configuration_tutorial.html}. 访问日期：2026年11月15日.
\bibitem{profiler} PyTorch. PyTorch 2.8性能剖析工具. \url{https://docs.pytorch.org/docs/2.8/profiler.html}. 访问日期：2026年11月15日.
\bibitem{pointvla} Yu等. Point What You Mean: Visually Grounded Instruction Policy. arXiv:2512.18933. \url{https://arxiv.org/abs/2512.18933}. 本报告仅借鉴显式视觉目标接口，本地适配不等于原论文完整复现。
\bibitem{cag} Fang等. When Vision Overrides Language: Evaluating and Mitigating Counterfactual Failures in VLAs. arXiv:2602.17659. \url{https://arxiv.org/abs/2602.17659v1}. 本地试验参考首版的动作引导思路，未复现其全部基准。
\bibitem{shortcut} Xing等. Shortcut Learning in Generalist Robot Policies: The Role of Dataset Diversity and Fragmentation. CoRL 2025, PMLR 305:3239--3266. \url{https://proceedings.mlr.press/v305/xing25a.html}.
\end{thebibliography}
